%% ============================================================
%%  macros.tex
%%  Định nghĩa lệnh viết tắt Toán học và các khối TikZ dùng chung
%% ============================================================

%% ---------- 1. Lệnh viết tắt Toán học ----------
\newcommand{\dx}{\,\mathrm{d}x}
\newcommand{\dy}{\,\mathrm{d}y}
\newcommand{\dt}{\,\mathrm{d}t}
\newcommand{\dif}[1]{\,\mathrm{d}#1}
\newcommand{\Sxoay}{V_{\text{xoay}}}
\newcommand{\ababs}[1]{\left| #1 \right|}
\newcommand{\ketqua}[1]{\[\boxed{#1}\]}

%% ---------- 2. Khối TikZ: minh họa diện tích 1 đường cong & trục Ox ----------
% Tham số: hàm số biểu diễn bằng \fx, khoảng [a,b]
\newcommand{\hinhdientich}[5]{%
  % #1: định nghĩa hàm (declare function), #2: a, #3: b, #4: xmin/xmax, #5: ymax
  \begin{tikzpicture}[scale=0.9, declare function={f(\x)=#1;}]
    \begin{axis}[
        axis lines=middle,
        xlabel={$x$}, ylabel={$y$},
        xmin=-1, xmax=#4, ymin=-0.5, ymax=#5,
        xtick={#2,#3}, ytick=\empty,
        width=8cm, height=6cm,
        axis line style={-Latex},
        samples=100,
      ]
      \addplot [domain=#2:#3, teal, fill=teal, fill opacity=0.25, thick] {f(x)} \closedcycle;
      \addplot [domain=-0.5:#4, navy, thick, samples=100] {f(x)};
      \draw[dashed, navylight] (axis cs:#2,0) -- (axis cs:#2,{f(#2)});
      \draw[dashed, navylight] (axis cs:#3,0) -- (axis cs:#3,{f(#3)});
    \end{axis}
  \end{tikzpicture}%
}

%% ---------- 3. Khối TikZ: minh họa diện tích giữa 2 đường cong ----------
\newcommand{\hinhhaihdc}[6]{%
  % #1: f(x), #2: g(x), #3: a, #4: b, #5: xmax, #6: ymax
  \begin{tikzpicture}[scale=0.9, declare function={f(\x)=#1; g(\x)=#2;}]
    \begin{axis}[
        axis lines=middle,
        xlabel={$x$}, ylabel={$y$},
        xmin=-1, xmax=#5, ymin=-0.5, ymax=#6,
        xtick={#3,#4}, ytick=\empty,
        width=8cm, height=6cm,
        axis line style={-Latex},
        samples=100,
      ]
      \addplot [domain=#3:#4, amber, fill=amber, fill opacity=0.25, thick] {f(x)} \closedcycle;
      \addplot [domain=#3:#4, fill=white, opacity=1] {g(x)} \closedcycle;
      \addplot [domain=-0.5:#5, navy, thick] {f(x)} node[right] {$y=f(x)$};
      \addplot [domain=-0.5:#5, teal, thick] {g(x)} node[right] {$y=g(x)$};
    \end{axis}
  \end{tikzpicture}%
}

%% ---------- 4. Khối TikZ tổng quát: sơ đồ tư duy nhánh (dùng cho MOD-03) ----------
\tikzset{
  khoinode/.style={rectangle, rounded corners, draw=navy, thick, fill=tealbg,
    text width=3.2cm, align=center, minimum height=1cm, font=\small},
  khoimoi/.style={-Latex, navy, thick}
}

%% ============================================================
%%  BỔ SUNG v2.0: Hộp cho chu trình tư duy "Khám phá - Gợi mở - Ý tưởng"
%% ============================================================

%% ---------- 16. Bảng màu bổ sung ----------
\definecolor{violet}{HTML}{6B4E9E}
\definecolor{violetbg}{HTML}{F1ECF9}
\definecolor{ideaframe}{HTML}{C6900A}
\definecolor{ideabg}{HTML}{FFF8E7}

%% ---------- 18. Hộp Câu hỏi gợi mở ----------
% Viền nét đứt - đóng vai trò như lời gợi ý từng bước để học sinh
% tự suy nghĩ TRƯỚC khi đọc phần phân tích ý tưởng.
\newtcolorbox{hopgoimo}[1][Quan sát \& Gợi mở]{
  enhanced, breakable,
  colback=white, colframe=navylight,
  frame style={dashed, line width=1pt},
  boxrule=1pt, arc=1mm, top=2mm, bottom=2mm,
  fonttitle=\bfseries, coltitle=navy,
  title={\faSearch\ \textbf{#1}}
}

\endinput

